\documentclass[11pt,a4paper]{article}
\usepackage[margin=1in]{geometry}
\usepackage{hyperref}
\usepackage{enumitem}
\usepackage{graphicx}
\usepackage{xcolor}


% Variable: Lab Instructor
\makeatletter
\newcommand{\BvrSetLabInstructor}[1]{%
  \gdef\Bvr@LabInstructor{#1}}
% default
\newcommand{\Bvr@LabInstructor}{Dr. Paramveer %
 Kaur}

% public accessor
\newcommand{\BvrLabInstructorValue}{\Bvr@LabInstructor}
\makeatother

% Add subtitle
\usepackage{titling}

% Set title bold
\pretitle{\begin{center}\bfseries\fontsize{18bp}{18bp}\selectfont}

\makeatletter
\newcommand{\BvrSubtitle}[1]{%
  \posttitle{%
    \par\end{center}
    \begin{center}\large#1\end{center}
    \vskip 0.5em}%
}
\makeatother

% Hints in section titles
\newcommand{\BvrHint}[1]{%
  {\normalfont\normalsize\itshape\hfill #1}}

% -----------------------------------------------------

\title{Distributed Performance Testing Platform}

\BvrSetLabInstructor{Dr. Paramveer Kaur
}

\BvrSubtitle{%
  Project Proposal for UCS503P  \\
  Submitted to: \BvrLabInstructorValue \\ \bfseries
  Thapar Institute of Engineering and Technology%
}
\author{
  Jatin Bajaj - \texttt{1024170389} \\[0.3em]
  Krish Sahni - \texttt{1024170390} \\[0.3em]
  Rishabh Dutt - \texttt{10240170393}
}


\date{\today}

\begin{document}
\maketitle

% -----------------------------------------------------
\section{Higher-order goal%
  \BvrHint{\ldots primary objective}}

This project aims to build a scalable, backend-focused distributed
performance testing platform capable of generating and analyzing
high volumes of concurrent HTTP traffic against APIs and web
applications, including applications running locally on a developer's
machine. It focuses primarily on distributed backend architecture,
scalable remote request generation, secure connectivity to local
applications, real-time telemetry processing, and automated
performance analysis.

The objective is not merely to generate load and display graphs, but
to provide an engineering platform capable of determining system
capacity, identifying performance bottlenecks, detecting breaking
points, comparing performance across different application versions,
and enabling users with limited hardware to use remote cloud-based
load generators for testing locally hosted applications.

% -----------------------------------------------------
\section{Time-to-value%
  \BvrHint{\ldots primary heuristic}}

\begin{itemize}[leftmargin=0.7cm]

    \item The initial iteration will focus on implementing the core
    controller-worker architecture and generating concurrent HTTP
    traffic against configurable API endpoints.

    \item Subsequent iterations will introduce distributed cloud
    workers, a lightweight local load-testing agent, secure remote
    connectivity, telemetry streaming, metric aggregation,
    real-time dashboards, adaptive stress testing, and performance
    regression testing.

    \item Development will prioritize measurable backend functionality
    over visual complexity, with each iteration producing a testable
    performance-testing capability.

\end{itemize}

% -----------------------------------------------------
\section{Problem Statement}

Existing performance-testing tools can generate large volumes of
requests, but a developer running a local web server on a low-resource
machine may not have sufficient CPU, memory, or network capacity to
generate large workloads from the same machine. In addition,
remotely generated traffic cannot directly access a developer's
\texttt{localhost} without a secure communication mechanism.

The proposed platform addresses the following challenges:

\begin{itemize}[leftmargin=0.7cm]

    \item \textbf{Limited scalability:} A single load generator may
    become a bottleneck when thousands of concurrent requests are
    required.

    \item \textbf{Distributed load generation:} Large workloads need
    to be distributed across multiple cloud machines or worker
    processes so that the load-generation infrastructure does not
    become the bottleneck.

    \item \textbf{Limited client hardware:} Developers and students
    may have low-resource machines that cannot efficiently generate
    thousands of concurrent requests while also running the
    application under test.

    \item \textbf{Local application accessibility:} Cloud-based workers
    cannot directly access a developer's \texttt{localhost}; a
    lightweight local agent and secure outbound communication channel
    are required to route remote test traffic to the local application.

    \item \textbf{Limited real-time analysis:} Raw request results need
    to be aggregated into useful metrics such as throughput,
    latency percentiles, and error rates.

    \item \textbf{Manual breaking-point detection:} Developers often
    need to manually increase traffic and inspect results to determine
    when a system becomes unstable.

    \item \textbf{Performance regressions:} Performance degradation
    between application versions may not be immediately visible
    without systematic comparison of historical test results.

\end{itemize}

% -----------------------------------------------------
\section{Proposed Solution%
  \BvrHint{\ldots software proposition}}

\subsection{Overview}

Build a backend-focused distributed load-testing platform consisting
of a central controller, multiple cloud-based load-generator workers,
a lightweight local load-testing agent, a secure communication
gateway, a metrics aggregation pipeline, persistent test storage,
and a real-time monitoring and visualization interface.

The local agent enables the platform to test applications that are
running on a developer's \texttt{localhost}. The agent establishes a
secure outbound connection to the platform and forwards test traffic
from remote workers to the locally hosted application. This allows
the load-generation workload to be executed remotely while the
developer's machine primarily runs the application under test.

The platform can therefore allow a developer with a low-resource
machine to run their application locally while using cloud-based
workers to generate thousands of concurrent requests. The collected
telemetry will be processed and presented through real-time graphs
and visualizations, allowing users to observe application performance
while the test is running.

The platform will provide:

\begin{itemize}[leftmargin=0.7cm]

    \item \textbf{Central controller:} Creates, configures, schedules,
    starts, and stops performance tests.

    \item \textbf{Load-generator workers:} Cloud-based, horizontally
    scalable workers that generate HTTP traffic concurrently and
    execute assigned workloads.

    \item \textbf{Local load-testing agent:} A lightweight process
    running on the user's machine that securely connects to the
    platform and forwards remote test traffic to applications running
    on \texttt{localhost}.

    \item \textbf{Communication gateway:} Maintains secure connections
    between cloud workers and local agents without exposing the
    developer's local application directly to the public Internet.

    \item \textbf{Metrics aggregation pipeline:} Processes high-volume
    telemetry generated by workers and calculates performance metrics
    such as throughput, requests per second, latency percentiles,
    error rates, timeouts, response codes, and concurrency.

    \item \textbf{Historical storage:} Stores test configurations,
    execution results, aggregated metrics, and performance history.

    \item \textbf{Real-time monitoring and visualization:} Displays
    live performance graphs for metrics such as requests per second,
    throughput, P50/P95/P99 latency, error rate, concurrency,
    response-code distribution, and SLO violations during test
    execution.

\end{itemize}

\subsection{Core workflow}

\begin{enumerate}[leftmargin=0.7cm]

    \item The user creates a performance test by configuring the target
    endpoint, virtual users, request rate, ramp-up period, duration,
    authentication, headers, payloads, and scenarios.

    \item If the target application is locally hosted, the user starts
    the lightweight local load-testing agent and registers the local
    target endpoint, such as \texttt{localhost:8080}.

    \item The local agent establishes a secure outbound connection to
    the platform so that remote workers can reach the local
    application through the controlled communication channel.

    \item The central controller validates the configuration and
    distributes the workload across available cloud load-generator
    workers.

    \item Workers generate concurrent HTTP requests and route them
    through the communication layer to the local agent when the
    target is locally hosted, or directly to a remotely accessible
    target when appropriate.

    \item Workers collect request-level telemetry and asynchronously
    stream it to the metrics aggregation pipeline.

    \item The aggregation layer calculates throughput, requests per
    second, latency percentiles, error rate, timeouts, response codes,
    and concurrent requests.

    \item Aggregated metrics are persisted and simultaneously streamed
    to the dashboard for real-time monitoring and visualization.

    \item After execution, the platform evaluates configured SLOs and
    determines whether the application remained within acceptable
    performance limits.

\end{enumerate}

% -----------------------------------------------------
\section{Performance Metrics%
  \BvrHint{\ldots measurable output}}

The platform will calculate and expose the following metrics:

\begin{itemize}[leftmargin=0.7cm]

    \item \textbf{Throughput:} Total number of successfully processed
    requests during a test.

    \item \textbf{Requests per second (RPS):} Number of requests
    generated or completed per second.

    \item \textbf{Latency percentiles:} P50, P95, and P99 latency
    values representing typical, high, and worst-case response
    behavior.

    \item \textbf{Error rate:} Percentage of requests resulting in
    failures.

    \item \textbf{Timeouts:} Number and percentage of requests that
    exceeded the configured timeout.

    \item \textbf{Response codes:} Distribution of HTTP status codes
    such as 2xx, 4xx, and 5xx responses.

    \item \textbf{Concurrency:} Number of active requests or virtual
    users during different stages of the test.

\end{itemize}

% -----------------------------------------------------
\section{Adaptive Stress Testing%
  \BvrHint{\ldots capacity discovery}}

Instead of using only a fixed load, the platform will support
adaptive stress testing.

The system will progressively increase traffic according to a
configured strategy and continuously evaluate application performance
against predefined Service Level Objectives (SLOs).

Example SLOs include:

\begin{itemize}[leftmargin=0.7cm]

    \item P95 latency must remain below a configured threshold.

    \item P99 latency must remain below a configured threshold.

    \item Error rate must remain below a configured percentage.

    \item Request timeout rate must remain below a configured limit.

    \item Minimum throughput must remain above a configured RPS
    threshold.

\end{itemize}

The \textbf{maximum sustainable throughput} will be identified as the
highest load at which the application continues to satisfy the
configured SLOs.

The \textbf{breaking point} will be identified when increasing load
causes the application to violate one or more critical SLOs, such as
a significant increase in latency, error rate, or timeouts.

The system can progressively narrow the load range around the first
SLO violation to estimate the capacity boundary more precisely.

This allows the system to automatically determine the approximate
capacity boundary of the tested application rather than requiring
manual inspection.

% -----------------------------------------------------
\section{Performance Regression Testing%
  \BvrHint{\ldots version comparison}}

The platform will support performance regression testing by storing
historical test results and associating them with application
versions or builds.

A new performance test can be compared against a previous baseline
to identify degradation in:

\begin{itemize}[leftmargin=0.7cm]

    \item P50/P95/P99 latency.

    \item Throughput and requests per second.

    \item Error and timeout rates.

    \item Maximum sustainable throughput.

    \item Resource or concurrency behavior when such measurements are
    available.

\end{itemize}

A performance regression occurs when a newer application version
performs significantly worse than the established baseline under
comparable test conditions.

This functionality also makes the platform suitable for integration
into a CI/CD pipeline, where performance tests can be executed
automatically for application builds.

% -----------------------------------------------------
\section{System Architecture%
  \BvrHint{\ldots backend engineering focus}}

The proposed platform will follow a distributed architecture in which
load generation is separated from the system under test. Cloud-based
workers perform the computationally intensive request generation,
while a lightweight local agent provides controlled access to
applications running on a developer's \texttt{localhost}.

The high-level architecture is:

\begin{center}
\texttt{Controller}
$\rightarrow$
\texttt{Cloud Workers}
$\rightarrow$
\texttt{Secure Gateway}
$\rightarrow$
\texttt{Local Agent}
$\rightarrow$
\texttt{localhost Application}
\end{center}

For remotely accessible applications, workers may communicate with
the target directly without using the local agent.

\begin{itemize}[leftmargin=0.7cm]

    \item \textbf{Controller:} Maintains test configuration, manages
    test lifecycle, schedules workloads, and coordinates workers.

    \item \textbf{Load Generators:} Horizontally scalable cloud
    workers that generate concurrent HTTP traffic without placing the
    primary load-generation cost on the user's local machine.

    \item \textbf{Local Agent:} A lightweight client running beside a
    locally hosted application and maintaining a secure outbound
    connection to the platform.

    \item \textbf{Secure Communication Gateway:} Routes requests
    between cloud workers and registered local agents while
    controlling authentication, authorization, connection state, and
    request routing.

    \item \textbf{Telemetry Pipeline:} Receives high-volume execution
    data asynchronously from workers.

    \item \textbf{Metrics Aggregator:} Converts raw telemetry into
    aggregated performance metrics and percentile statistics.

    \item \textbf{Persistent Storage:} Stores test definitions,
    execution metadata, historical metrics, and regression baselines.

    \item \textbf{WebSocket Layer:} Provides real-time performance
    updates to connected dashboard clients.

\end{itemize}

The architecture is designed so that additional cloud
load-generator workers can be added horizontally when higher traffic
generation capacity is required. This separates the capacity of the
load-generation infrastructure from the hardware available to the
developer running the application under test.

A developer with a low-resource machine can therefore run a local
application and use multiple remote workers to generate a large
concurrent workload. The local agent performs request forwarding and
communication management rather than acting as the primary load
generator.

% -----------------------------------------------------
\section{Solution Approach%
  \BvrHint{\ldots engineering focus}}

The project will focus on backend engineering challenges rather than
AI/ML-based prediction.

\subsection{Concurrent Request Generation}

Workers will use asynchronous or non-blocking networking techniques
to efficiently generate a large number of concurrent HTTP requests
without requiring one dedicated blocking thread for every request.

\subsection{Distributed Coordination}

The controller will divide workload among cloud workers and maintain
the state of active performance tests. Workers will report their
execution status and telemetry back to the controller or telemetry
pipeline. Worker registration, health checks, heartbeats, and
workload assignment will allow the controller to manage a changing
set of load-generation nodes.

\subsection{Local Application Connectivity}

A lightweight local agent will maintain a secure outbound connection
to the platform. When a test targets \texttt{localhost}, the platform
will route requests through this connection to the locally running
application.

The design avoids requiring the user's application to be directly
exposed to the public Internet and prevents the local machine from
being responsible for generating the complete workload.

\subsection{Rate Limiting}

Request generation will support configurable request rates to prevent
uncontrolled traffic generation and to implement controlled load
profiles.

\subsection{Telemetry Streaming}

Workers will stream high-volume execution data rather than waiting
until the complete test finishes. This enables real-time aggregation
and monitoring.

\subsection{Fault Tolerance}

The platform will account for worker failures and communication
issues so that individual worker failures do not unnecessarily
terminate the complete testing infrastructure.

% -----------------------------------------------------
\section{Real-time Monitoring%
  \BvrHint{\ldots live system visibility}}

During a running test, aggregated metrics will be transmitted to the
dashboard through WebSockets.

The dashboard will provide real-time visibility into:

\begin{itemize}[leftmargin=0.7cm]

    \item Current requests per second.

    \item Active concurrent requests.

    \item P50, P95, and P99 latency.

    \item Error and timeout rates.

    \item HTTP response-code distribution.

    \item Current load level during adaptive stress testing.

    \item SLO violations and detected performance limits.

\end{itemize}

The visualization layer will present these metrics through
real-time graphs and charts, allowing users to observe changes in
latency, throughput, concurrency, error rates, and load levels during
test execution.

This avoids the need to wait until the test finishes before
understanding the application's current behavior.

% -----------------------------------------------------
\section{CI/CD Integration%
  \BvrHint{\ldots automated regression detection}}

The platform can be integrated into CI/CD pipelines so that
performance tests become an automated quality gate.

A typical pipeline can execute the following workflow:

\begin{enumerate}[leftmargin=0.7cm]

    \item Build a new application version.

    \item Deploy the application to a test environment.

    \item Trigger the performance-testing platform through an API.

    \item Execute a predefined workload.

    \item Collect and aggregate performance metrics.

    \item Compare results against the stored baseline.

    \item Mark the pipeline as successful or failed according to
    configured performance SLOs.

\end{enumerate}

This enables performance degradation to be detected before a build is
released to production.

% -----------------------------------------------------
\section{Evaluation Criteria%
  \BvrHint{\ldots measurable and attributable}}

The project will be evaluated using measurable technical metrics.

\subsection{Primary Evaluation Metrics}

\begin{itemize}[leftmargin=0.7cm]

    \item \textbf{Load generation capacity:} Maximum requests per
    second generated by the platform.

    \item \textbf{Concurrency:} Maximum number of concurrent virtual
    users and requests supported.

    \item \textbf{Telemetry throughput:} Number of execution events
    processed by the aggregation pipeline per second.

    \item \textbf{Metric accuracy:} Correctness of calculated
    throughput, latency percentiles, error rates, and response-code
    distributions.

\end{itemize}

\subsection{Secondary Evaluation Metrics}

\begin{itemize}[leftmargin=0.7cm]

    \item Controller-to-worker coordination latency.

    \item Worker failure recovery behavior.

    \item Real-time dashboard update latency.

    \item Storage scalability for historical test results.

    \item Accuracy of maximum sustainable throughput and breaking-point
    detection.

    \item Local-agent request forwarding latency and reliability.

    \item Scalability achieved by adding additional cloud workers.

\end{itemize}

% -----------------------------------------------------
\section{Scalability%
  \BvrHint{\ldots horizontal growth}}

The system is designed for horizontal scalability.

\begin{enumerate}[leftmargin=0.7cm]

    \item \textbf{Load-generator scalability:} Additional cloud
    workers can be added to increase traffic-generation capacity
    without requiring upgrades to the user's local machine.

    \item \textbf{Remote testing of local applications:} A lightweight
    agent allows cloud workers to test applications bound to
    \texttt{localhost}, making the platform usable for developers and
    students with limited hardware.

    \item \textbf{Aggregation scalability:} Telemetry processing can
    be distributed across multiple aggregation instances when event
    volume increases.

    \item \textbf{Stateless services:} Controller and API services can
    be designed to remain stateless where possible, allowing multiple
    instances to run concurrently.

    \item \textbf{Storage scalability:} Historical performance data
    can be indexed and partitioned according to test, timestamp, and
    application version.

    \item \textbf{Asynchronous processing:} Streaming and asynchronous
    processing prevents high-volume telemetry from blocking request
    generation.

\end{enumerate}

% -----------------------------------------------------
\section{Project Scope and Deliverables%
  \BvrHint{\ldots iteration-friendly}}

\subsection{Initial Deliverable%
  \BvrHint{\ldots first iteration}}

\begin{itemize}[leftmargin=0.7cm]

    \item Central controller for creating and managing tests.

    \item Single cloud load-generator worker capable of generating
    configurable concurrent HTTP requests.

    \item Basic local load-testing agent capable of establishing a
    secure connection to the platform and forwarding requests to a
    locally hosted application.

    \item Configurable virtual users, request rate, duration,
    ramp-up/ramp-down, headers, authentication, and payloads.

    \item Collection of request-level performance metrics.

    \item Basic metrics aggregation for throughput, latency,
    response codes, and errors.

    \item Persistent storage of test configurations and results.

\end{itemize}

\subsection{Subsequent Deliverables}

\begin{itemize}[leftmargin=0.7cm]

    \item Multiple distributed cloud load-generator workers.

    \item Local load-testing agent with secure connection management.

    \item Communication gateway for routing remote traffic to local
    applications.

    \item Streaming telemetry pipeline.

    \item Real-time WebSocket dashboard.

    \item Real-time performance graphs and visualization.

    \item P50/P95/P99 latency analysis.

    \item Adaptive stress testing and breaking-point detection.

    \item Configurable SLO-based performance evaluation.

    \item Historical performance comparison.

    \item Performance regression testing.

    \item CI/CD pipeline integration.

    \item Worker failure handling and fault tolerance.

\end{itemize}

% -----------------------------------------------------
\section{Risks and Mitigations}

\begin{itemize}[leftmargin=0.7cm]

    \item \textbf{Load-generator bottleneck:} Use asynchronous networking
    and distribute traffic across multiple cloud workers.

    \item \textbf{Local machine bottleneck:} Move request generation
    to remote workers so that the user's machine primarily processes
    the application under test rather than generating the complete
    workload.

    \item \textbf{Telemetry overload:} Use batching, asynchronous
    streaming, and aggregation instead of storing every event
    synchronously.

    \item \textbf{Incorrect percentile calculations:} Use validated
    percentile algorithms and test metric calculations against known
    datasets.

    \item \textbf{Worker failure:} Track worker heartbeats and isolate
    failed workers from the remaining test execution.

    \item \textbf{Secure local connectivity:} Authenticate local
    agents, use secure outbound connections, restrict request routing,
    and avoid exposing arbitrary local ports to remote users.

    \item \textbf{Local network bottleneck:} Clearly distinguish
    remote-to-local testing from cloud-to-cloud testing because
    bandwidth and network latency between the user's machine and the
    cloud can affect test results.

    \item \textbf{False breaking-point detection:} Evaluate multiple
    metrics and require configurable SLO violations rather than
    relying on a single latency measurement.

    \item \textbf{Storage growth:} Store aggregated metrics for
    long-term analysis while limiting unnecessary raw telemetry
    retention.

\end{itemize}

% -----------------------------------------------------
\section{Summary} This project proposes a backend-focused distributed performance testing platform capable of generating large-scale HTTP workloads using a central controller and horizontally scalable cloud load-generator workers. A lightweight local agent will additionally allow users to test applications running on \texttt{localhost} without requiring their local machine to generate the complete workload. Unlike a basic load-testing tool that only generates requests and plots graphs, the proposed system focuses on distributed workload orchestration, secure remote-to-local request routing, asynchronous networking, telemetry streaming, real-time metric aggregation, performance visualization, adaptive stress testing, SLO-based capacity analysis, breaking-point detection, and performance regression testing. The visualization layer will transform aggregated telemetry into real-time performance graphs, latency percentile trends, throughput and concurrency charts, response-code distributions, SLO indicators, capacity visualizations, and historical comparisons. The platform will provide a practical environment for studying and implementing backend engineering concepts including high concurrency, distributed systems, rate limiting, asynchronous processing, secure tunneling, worker coordination, fault tolerance, scheduling, real-time communication, performance benchmarking, CI/CD quality gates, performance visualization, and scalable data storage.

% -----------------------------------------------------

\end{document}